% Capítulo 2, versão 7 (18 nov. 2026). Base: revisão 6, com as edições do autor de 17 e 18 nov. 2026.
% Comentários "% V7" marcam o que mudou nesta versão.
% Capítulo 2, versão 6 (16 nov. 2026).
% "% V6" marca o que mudou em relação à versão 5. Os detalhes da solução foram reduzidos aqui e ficam no Capítulo 5.
% As descrições dos nove trabalhos relacionados foram conferidas nos PDFs da pasta TCC-fonte em 16 nov. 2026.

\chapter{Referencial Teórico}
\label{cap-referencial-teorico}

% V7: parágrafo de abertura escrito pela orientadora e ajustado pelo autor. As remissões usam "Seção" com maiúscula, como no restante do trabalho.
Com foco em segurança de software e ênfase em auditoria, o capítulo começa apresentando a auditoria de repositórios de software (seção~\ref{rt-2-1}), cobrindo, dentre outros tópicos, aspectos de qualidade e DevSecOps. Na sequência, e como uma reflexão sobre as práticas de mercado nessa área, há menção à diversidade das automações, às plataformas \textit{low-code/no-code} e à \textit{Shadow AI} (seção~\ref{rt-2-2}), entendida como o uso de diferentes ferramentas de inteligência artificial por funcionários de uma empresa sem o conhecimento, a aprovação ou o controle dessa empresa. O paradigma de Sistemas Multiagentes é abordado na seção~\ref{rt-2-3}, considerando, dentre outros pontos, a comunicação e a cooperação entre agentes (seção~\ref{rt-2-4}). Há um olhar ainda para governança e apoio à decisão (seção~\ref{rt-2-5}) e para os trabalhos relacionados, comparados à proposta deste trabalho (seção~\ref{rt-2-6}). Por fim, a seção~\ref{rt-resumo} resume o capítulo.

\section{Auditoria de Repositórios de Software}
\label{rt-2-1}

Uma auditoria de software é uma avaliação independente da conformidade de produtos e processos de software com os critérios aplicáveis, como normas, planos e procedimentos \cite{ieee1028}. Um repositório reúne o código de uma solução e os artefatos que o acompanham, como dependências, arquivos de configuração e documentação. Com base na definição da norma, este trabalho chama de auditoria de repositório a avaliação de uma versão identificada desses artefatos em relação a critérios de qualidade e de segurança.

% V6: parágrafo novo. Evita apresentar a separação entre agentes como independência de auditoria.
% V7: conferido na ISO 19011:2026 (4. ed.), item 4.6, p. 6, na amostra pública enviada pelo autor.
A independência mencionada na definição refere-se à relação entre quem avalia e o que é avaliado. A ISO 19011 a trata como um princípio da auditoria, base da imparcialidade e da objetividade das conclusões, e recomenda que os auditores sejam independentes da atividade auditada sempre que possível \cite[item 4.6]{iso19011}. Por isso, a solução proposta neste trabalho é entendida como apoio à auditoria, que organiza achados e evidências para o exame de um responsável humano, e não como uma auditoria independente.

% V4: o termo "achado candidato" saiu. Ele não estava na lista de termos da seção e não tinha fonte.
% V7: redação do autor ("dois conceitos"), com os itens da ISO 19011:2026 conferidos no texto da norma. A frase sobre a analogia foi mantida, pois a norma trata de sistemas de gestão.
% V4: Oliveira (2026), indicada pela orientadora, foi conferida (seção 4.6, p. 43-44 da monografia; a seção 4.1, p. 38, também define o termo).
Dois conceitos da literatura de auditoria descrevem o que esse exame produz, que são as evidências e os achados. Evidências são registros, declarações de fatos ou outras informações pertinentes aos critérios da auditoria e que podem ser verificadas \cite[item 3.10]{iso19011}. Achados são os resultados da avaliação dessas evidências em relação aos critérios \cite[item 3.11]{iso19011}. A norma trata de auditorias de sistemas de gestão, e os dois conceitos são adotados aqui por analogia.

Em segurança de software, o termo achado é empregado de forma ampla para qualquer ocorrência reportada automaticamente por uma ferramenta de análise \cite{oliveira2026}. Este trabalho adota esse uso. Assim, quando uma ferramenta aplica uma regra ao código e aponta uma ocorrência, o resultado é um achado, e o relatório da ferramenta e o trecho de código apontado são as evidências que permitem conferi-lo. Um achado não equivale, por si só, a uma vulnerabilidade confirmada, pois ainda exige revisão \cite{oliveira2026}.

% V7: item 3.12 da ISO 19011:2026, conferido no texto da norma. A abertura mudou porque o parágrafo anterior passou a apresentar dois conceitos.
A norma define ainda a conclusão da auditoria, isto é, o resultado obtido depois de considerados os objetivos da auditoria e todos os achados \cite[item 3.12]{iso19011}. % V5: a introdução final define o parecer como uma recomendação.
Neste trabalho, o documento que registra essa conclusão é chamado de parecer. Ele traz uma recomendação sobre a versão examinada e indica as evidências que a sustentam.

As subseções seguintes apresentam as atividades de uma auditoria de repositório. A última relaciona essas atividades à prática DevSecOps e ao uso que este trabalho faz delas.

\subsection{Verificação de qualidade}
\label{rt-2-1-1}

A primeira atividade examina a qualidade do produto. A ISO/IEC 25010 define um modelo de qualidade composto por nove características, divididas em subcaracterísticas, que apoia a especificação de requisitos e a avaliação da qualidade ao longo do ciclo de vida \cite{iso25010}. Como o modelo reúne várias características, uma medida isolada não descreve a qualidade de um produto.

Parte dessas características é examinada por métricas. A complexidade ciclomática, proposta por \citeonline{mccabe1976}, mede o número de caminhos linearmente independentes no grafo de fluxo de controle do programa. No SonarQube, o contador aumenta sempre que o fluxo de controle de uma função se divide, e o cálculo varia ligeiramente conforme a linguagem \cite{sonarMetrics99}. Por isso, uma medida deve ser lida na unidade em que foi calculada, e o valor agregado de um arquivo não descreve uma função em particular.

\subsection{Análise de código e de componentes}
\label{rt-2-1-2}

A segunda atividade procura problemas de segurança no código. O \textit{Secure Software Development Framework} (SSDF), do NIST, recomenda revisar ou analisar o código para identificar vulnerabilidades e verificar a conformidade com os requisitos de segurança. Recomenda também registrar e triar os problemas encontrados \cite{nistSsdf}. Essa análise pode ser estática ou dinâmica. A análise estática raciocina sobre o código como está escrito, sem depender de sua execução, e a análise dinâmica se baseia na execução do programa \cite{blackIr8397}.

Os achados da análise estática precisam ser examinados antes de sustentar uma conclusão. Em entrevistas com 20 desenvolvedores, \citeonline{johnson2013} identificaram os falsos positivos e a forma de apresentação dos alertas entre as barreiras ao uso dessas ferramentas. Um alerta de ferramenta, portanto, indica um ponto a examinar e não constitui, por si só, um problema confirmado.

% V5: a introdução final prevê confrontar os achados com boas práticas da comunidade de segurança. As duas fontes foram conferidas em 16 nov. 2026.
Os achados de segurança costumam ser classificados por referências mantidas pela comunidade. A \textit{Common Weakness Enumeration} (CWE) é uma lista de fraquezas comuns de software e de hardware. Uma fraqueza é uma condição que, em certas circunstâncias, pode contribuir para a introdução de vulnerabilidades \cite{cweAbout}. O OWASP Top 10, em sua edição de 2025, agrupa 248 dessas fraquezas em dez categorias de risco para aplicações \cite{owaspTop10}. Neste trabalho, essas referências servem para confrontar os achados documentados, isto é, para verificar se a classificação registrada no parecer corresponde à fonte que a ferramenta indica.

A terceira atividade trata dos componentes de terceiros. As ferramentas de análise de composição de software (\textit{Software Composition Analysis}, SCA) ajudam a identificar as bibliotecas e os pacotes que o software utiliza e a apontar os que estão desatualizados ou têm vulnerabilidades conhecidas. Como uma nova vulnerabilidade pode ser divulgada a qualquer momento, os componentes precisam ser acompanhados continuamente \cite{blackIr8397}. A edição de 2025 do OWASP Top 10 passou a tratar o tema em uma categoria própria, dedicada às falhas na cadeia de suprimentos de software \cite{owaspTop10}.

\subsection{DevSecOps e interpretação dos resultados}
\label{rt-2-1-3}

Essas atividades costumam ser reunidas na prática DevSecOps, isto é, na integração de princípios e práticas de segurança ao DevOps \cite{rajapakse2022}. Nessa prática, as etapas de construção e entrega do software são organizadas em esteiras de integração e entrega contínuas (CI/CD) \cite{chandramouli2024}.

A literatura registra dificuldades nessa integração. Em uma revisão sistemática de 54 estudos, \citeonline{rajapakse2022} observam que a segurança é tradicionalmente tratada como requisito não funcional, verificada sobre uma versão já pronta, e que atividades como a revisão de código exigem esforço manual substancial. Os autores também relatam o tempo consumido pela avaliação manual dos falsos positivos das ferramentas.

% V4: frase ajustada ao que o artigo registra (questão Q13, sobre desafios de desenvolvimento seguro; 28,4% e 14,8% das respostas).
Em uma pesquisa com 88 profissionais de organizações brasileiras, \citeonline{canedo2026} registram, entre os desafios do desenvolvimento seguro, as vulnerabilidades no código e em bibliotecas externas e a automação limitada das análises de segurança.

Um trabalho próximo desta proposta é o de \citeonline{zuev2025}, que reúnem os relatórios de análises estática e dinâmica em um módulo de apoio à decisão sobre a segurança do código. Os autores argumentam que, diante de grandes volumes de dados, o fator humano pode levar a erros. Relatam ainda um experimento com cerca de 100 varreduras, no qual o módulo tornou mais eficiente a análise das vulnerabilidades identificadas.

% V4: fonte indicada pela orientadora, conferida no texto do artigo. O trabalho é conceitual.
Outro trabalho próximo é o de \citeonline{carvalho2026}, que discutem o uso de agentes autônomos de inteligência artificial em esteiras DevSecOps. Os autores defendem que cada registro de vulnerabilidade conserve sua origem, a versão da ferramenta, o \textit{commit} e a evidência reproduzível. Defendem também que a aceitação de riscos permaneça com uma pessoa ou instância identificável. O trabalho é conceitual, e os próprios autores apontam a necessidade de validação empírica.

% V6: fechamento reescrito para explicitar a cadeia achado, evidência, contexto documental, conclusão do especialista e parecer.
Neste trabalho, a auditoria de repositório reúne essas atividades, e seus termos formam uma cadeia. As ferramentas produzem achados de qualidade, de segurança e de dependências, e as evidências permitem conferi-los. Falta, porém, saber o que cada achado significa para a organização. Na metodologia de avaliação de riscos da OWASP, o impacto para o negócio é, em última análise, o fator mais importante \cite{owaspRiskRating}. Por isso, a proposta acrescenta aos achados o contexto documental, isto é, as informações registradas sobre a finalidade e a criticidade do componente examinado.

Um exemplo ilustra essa relação. Quando uma ferramenta aponta uma vulnerabilidade em um componente, tem-se um achado, e o relatório da ferramenta e o trecho de código apontado são as evidências que permitem conferi-lo. Se os documentos do projeto informam que o componente trata dados pessoais, essa informação compõe o contexto documental do achado e ajuda a interpretá-lo, pois a mesma vulnerabilidade tem outro peso em um componente com essa finalidade. O parecer reúne esses elementos e apoia a decisão do responsável humano, sem substituí-la.

\section{Diversidade das Automações, \textit{Low-code/No-code} e \textit{Shadow AI}}
\label{rt-2-2}

O contexto documental importa ainda mais quando as automações têm origens diferentes. O desenvolvimento pode envolver profissionais de software e usuários de outras áreas da organização. As plataformas \textit{low-code/no-code} (LCNC) reduzem a necessidade de programação manual e ampliam a participação desses usuários, prática associada ao \textit{citizen development}. Essa adoção traz benefícios e desafios organizacionais e técnicos, incluindo questões de segurança e governança \cite{ajimati2025}.

Outro fator é a \textit{Shadow AI}, isto é, o uso de inteligência artificial fora das estruturas de governança aprovadas pela organização. Esse uso envolve questões de privacidade dos dados e de responsabilidade sobre a tecnologia \cite{silic2025}. No contexto brasileiro, profissionais ouvidos por \citeonline{canedo2026} relataram preocupação com o uso de IA fora da supervisão institucional.

O desenvolvimento por usuários de negócio e o uso não autorizado de IA são fenômenos distintos, e a ocorrência de cada um precisa ser identificada nos registros da organização. Neste trabalho, a presença dessas práticas no portfólio da Empresa A será caracterizada com base nos registros institucionais disponíveis. Qualquer que seja a origem de uma automação, interpretar seus achados exige conhecer a finalidade do componente, e essa necessidade motiva o agente de Requisitos e Documentação, apresentado na próxima seção.

\section{Sistemas Multiagentes}
\label{rt-2-3}

A auditoria de repositório envolve atividades distintas, e cada uma pede um tipo de conhecimento. Um sistema multiagente permite distribuir essas atividades entre participantes especializados. Esta seção define agente e sistema multiagente, apresenta um modelo usado nos trabalhos relacionados e delimita o papel do modelo de linguagem na proposta.

\subsection{Agentes de software e responsabilidades}
\label{rt-2-3-1}

Um agente de software pode ser caracterizado por autonomia, capacidade de interação social, reação às mudanças percebidas e iniciativa orientada a objetivos \cite{wooldridge1995}. Um sistema multiagente (SMA) reúne agentes que atuam e interagem em um ambiente para alcançar objetivos individuais ou coletivos \cite{wooldridge2009,lazarin2025}. Essas interações podem envolver cooperação, coordenação e negociação \cite{wooldridge2009}.

% V5: Wang et al. foi conferido na versão arXiv (2505.07176v2), aceita pelo IEEE TCCN. O artigo é a revisão "Internet of Agents".
% A grafia "semi-centralizada" segue a introdução final.
Os sistemas multiagentes são associados a cenários distribuídos. Em uma revisão sobre a interconexão de agentes, \citeonline{wang2026} observam que as implantações reais envolvem agentes distribuídos em muitos dispositivos e locais. Os autores descrevem arquiteturas híbridas, nas quais conexões P2P entre agentes convivem com agentes que concentram o registro, a descoberta e a intermediação. Eles associam os sistemas multiagentes tradicionais a uma orquestração centralizada ou semi-centralizada.

% V7: redação do autor. Moro et al. (2003) saiu deste parágrafo. Só o resumo foi lido, e a fonte não sustentava mais do que "promissora".
% CONFERIR: Flores (2005), seção 2.3, p. 17-18, sustenta os tipos de rede. Confira se a definição geral de P2P (primeira frase após a abertura) está na monografia e indique a página.
% Moro (2003) e Flores (2005) não são usados para afirmar preferência geral por controle semi-centralizado.
Essa distribuição pede uma arquitetura compatível com ela. Segundo \citeonline{flores2005}, em uma rede \textit{peer-to-peer} (P2P), os participantes têm papéis equivalentes e se comunicam diretamente, cada um podendo fornecer e consumir recursos. Quanto à organização, o autor classifica essas redes em descentralizadas e semi-centralizadas. Na organização descentralizada, não há ponto central, e todos os nós têm a mesma importância. Na semi-centralizada, os nós têm importâncias diferentes, e um nó central ou alguns nós concentram as informações de controle. A arquitetura híbrida descrita por \citeonline{wang2026} se aproxima desse segundo arranjo, pois combina conexões diretas entre agentes com agentes que concentram funções de registro e intermediação.

% V6: fechamento reduzido. Os papéis de cada agente ficam no Capítulo 5.
A distribuição de funções entre agentes também é tratada na literatura como medida de segurança. \citeonline{mosciatti2026} adotam a separação de deveres como princípio de projeto. Nenhum agente reúne o contexto completo da tarefa, todas as credenciais e o acesso irrestrito às ferramentas. Trata-se de uma separação de responsabilidades dentro do sistema, que limita o efeito de um comportamento indevido. Ela não equivale à independência de auditoria discutida na seção~\ref{rt-2-1}.

Nesta proposta, as atividades da auditoria são distribuídas entre agentes especialistas, que se consultam diretamente por P2P, e o controle é semi-centralizado no Agente Auditor, que consolida o parecer. A arquitetura é detalhada no Capítulo~\ref{cap-proposta}.

\subsection{Modelos de agentes}
\label{rt-2-3-2}

Parte dos trabalhos relacionados descreve seus agentes pelo modelo \textit{Belief-Desire-Intention} (BDI). Nesse modelo, crenças representam o conhecimento e as percepções do agente, desejos representam seus objetivos, e intenções representam os compromissos assumidos com os objetivos selecionados. Planos descrevem as maneiras de persegui-los \cite{raoGeorgeff1995,lazarin2025}.

\citeonline{yesil2026} aplicam agentes BDI ao controle de acesso a documentos. As verificações especializadas são distribuídas entre agentes, e a lógica de autorização é expressa em planos que podem ser inspecionados. Extração de informação, verificação de políticas e decisão final ficam, assim, a cargo de agentes diferentes.

Esta proposta aproveita dessa literatura a especialização dos agentes e a consolidação explícita dos resultados.

\subsection{Interpretação documental apoiada por modelos de linguagem}
\label{rt-2-3-3}

Um modelo de linguagem pode servir a um agente como recurso para interpretar textos. No trabalho de \citeonline{yesil2026}, modelos de linguagem extraem informações dos documentos, e a decisão de autorização fica em uma camada de regras explícitas. Interpretar o conteúdo e decidir são, portanto, funções separadas.

% V6: fechamento reduzido. Os campos das respostas e as regras de associação ficam na Seção 5.6.
A separação tem um limite que os próprios autores reconhecem em parte. Um documento é uma entrada externa, e seu conteúdo pode conter trechos que um modelo de linguagem interpreta como instruções. Esse risco é conhecido como injeção indireta de \textit{prompt}. Entre as medidas indicadas para reduzi-lo estão separar e identificar o conteúdo externo, validar o formato da saída por código determinístico e restringir os privilégios do modelo ao mínimo necessário \cite{owaspLlm01}.

Esta pesquisa segue a mesma separação. O modelo de linguagem é usado apenas na interpretação dos documentos do projeto, para fornecer o contexto documental de um achado, e a recomendação resulta de regras explícitas. Uma menção genérica a uma funcionalidade em um documento não indica quais arquivos a implementam, e uma criticidade declarada é diferente de uma criticidade inferida.

\section{Comunicação e Cooperação entre Agentes}
\label{rt-2-4}

Agentes especializados precisam trocar informações para concluir suas análises. Esta seção apresenta como as mensagens são representadas e define a comunicação P2P e a cooperação adotadas neste trabalho.

\subsection{Mensagens, finalidade e transporte}
\label{rt-2-4-1}

Para que uma mensagem seja interpretada, importam o conteúdo transmitido e a finalidade da comunicação. Linguagens de comunicação de agentes, como KQML e FIPA-ACL, representam essas finalidades. Elas permitem distinguir, por exemplo, o fornecimento de uma informação de uma consulta dirigida a outro agente \cite{finin1994,fipaACL2002,lazarin2025}.

A representação da mensagem e o seu transporte são aspectos distintos. \citeonline{lazarin2025} transportam mensagens KQML por protocolos de correio eletrônico e traduzem os elementos entre a linguagem dos agentes e o formato de envio. O canal, portanto, pode variar, desde que os participantes consigam interpretar a mensagem. Neste trabalho, o formato das mensagens e o protocolo de transporte são especificados nos Capítulos~\ref{cap-suporte-tecnologico} e~\ref{cap-proposta}.

\subsection{Interações P2P e reavaliação das conclusões}
\label{rt-2-4-2}

% V4: sem ponto e vírgula. O transporte fica para o Cap. 3.
Neste trabalho, a comunicação P2P é a troca direta de uma pergunta e de uma resposta entre agentes especialistas. Uma pergunta do agente de Segurança pode ser dirigida ao de Qualidade ou ao de Requisitos e Documentação, e a resposta retorna ao solicitante. % V7: conteúdo do autor (SPADE, XMPP e papéis do Orquestrador e do Auditor), sem o ponto e vírgula.
O termo P2P designa aqui a interação lógica entre os especialistas. O transporte previsto utiliza SPADE sobre XMPP, com um servidor de mensagens \cite{spadeDocs,rfc6120}, e por isso essa interação não implica uma infraestrutura de rede sem servidores. O Orquestrador organiza as tarefas, e o Auditor centraliza a consolidação das conclusões. O transporte é descrito no Capítulo~\ref{cap-suporte-tecnologico}.

A cooperação, a coordenação e a negociação fazem parte do contexto de interação dos sistemas multiagentes \cite{lazarin2025}. Aqui, a cooperação toma a forma de uma consulta. Um especialista pede uma informação que outro pode fornecer e examina sua pertinência antes de concluir a própria análise. A resposta pode sustentar a manutenção ou a revisão da conclusão.

O significado dessas trocas depende do contexto. Ao estudar bots da Wikipédia, \citeonline{geiger2017} mostraram que interações antes classificadas como conflitos eram, em muitos casos, atividades rotineiras e colaborativas. Por isso, as perguntas, as respostas e as conclusões serão registradas, o que permite identificar quais informações foram compartilhadas e como entraram na análise.

O termo negociação não é usado para essas consultas. Ele fica reservado a mecanismos com propostas, critérios e procedimentos de resolução de divergências, que não fazem parte desta proposta.

\section{Governança, Pareceres e Revisão Humana}
\label{rt-2-5}

As conclusões dos agentes só apoiam uma decisão se forem consolidadas por critérios conhecidos e registradas de modo verificável. Esta seção trata das políticas de governança que dão origem a esses critérios, do parecer que os aplica e dos registros que permitem explicá-lo e revisá-lo.

\subsection{Governança de segurança de software, políticas e regras}
\label{rt-2-5-1}

% V5: subseção ampliada para sustentar o Objetivo Específico 6 (insumos de Governança de Software em Segurança).
% As três fontes novas foram conferidas em 16 nov. 2026. Veja o banco de citações.
A governança de segurança define como uma organização estabelece, comunica e acompanha sua estratégia, suas expectativas e suas políticas de gestão de riscos. No \textit{Cybersecurity Framework} 2.0, do NIST, essa é a função Governar, que orienta as demais funções de acordo com a missão da organização. Entre seus resultados esperados está a compreensão dos requisitos legais, regulatórios e contratuais aplicáveis \cite{nistCsf2}.

No desenvolvimento de software, a governança se traduz em políticas e em verificações de conformidade. O modelo de maturidade OWASP SAMM reúne essas atividades na prática Política e Conformidade. Essa prática busca atender às exigências legais e regulatórias externas e manter padrões internos de segurança alinhados ao propósito do negócio. Um de seus caminhos de melhoria é transformar padrões e obrigações em requisitos das aplicações, o que facilita auditorias mais eficientes e automatizadas \cite{owaspSamm}.

Entre as exigências legais brasileiras está a Lei Geral de Proteção de Dados Pessoais (LGPD). O Art.~46 determina que os agentes de tratamento adotem medidas de segurança, técnicas e administrativas, para proteger os dados pessoais. Segundo o § 2º do mesmo artigo, essas medidas devem ser observadas ``desde a fase de concepção do produto ou do serviço até a sua execução'' \cite{lgpd2018}. Uma automação que trata dados pessoais está, portanto, sujeita a controles que uma auditoria de repositório pode ajudar a verificar.

Para que tais políticas orientem um sistema de agentes, elas precisam ser representadas de forma explícita. Segundo \citeonline{ait2025}, a governança de projetos de software trata de tornar explícitos os papéis, as responsabilidades e os procedimentos pelos quais as decisões são tomadas. Os autores propõem uma linguagem específica de domínio para representar participantes, papéis, escopos e procedimentos de governança de forma estruturada. Em sistemas multiagentes normativos, as regras podem representar obrigações, permissões e proibições. Processar essas normas é diferente de aplicar sanções, que exigem representar consequências e mecanismos próprios de execução \cite{yan2025}.

Nesta proposta, esses insumos de governança chegam ao Agente Auditor na forma de regras versionadas. Uma diretriz escrita em um documento corporativo não é, por si só, uma regra executável. Por isso, este trabalho distingue as regras demonstrativas, elaboradas apenas para exercitar a solução, das regras institucionais, que traduzem políticas da Empresa A, leis ou normas e indicam a fonte de que derivam.

\subsection{Pareceres e delimitação da decisão}
\label{rt-2-5-2}

O apoio à decisão organiza informações e critérios para que o responsável decida com base em material que pode ser examinado. No trabalho de \citeonline{zuev2025}, a coleta dos resultados, seu processamento e a recomendação ao usuário são etapas separadas, e varreduras malsucedidas ou em número insuficiente limitam o que se pode concluir.

% V6: as categorias de recomendação saíram daqui e ficam na Seção 5.7.
Neste trabalho, o parecer reúne os achados, as conclusões dos especialistas, o contexto documental e as regras aplicadas. O Agente Auditor consolida esses elementos por regras explícitas e registra uma recomendação centrada em segurança. O parecer também sinaliza pontos de atenção, isto é, achados que uma regra indica para exame prioritário, e lacunas, isto é, informações que uma regra exige e que não foram encontradas.

A recomendação não é uma decisão. Ela não aprova nem impede, de forma automática, a entrega de uma automação. A decisão cabe ao responsável humano, e o parecer identifica as situações que precisam ser encaminhadas a ele. 

\subsection{Explicabilidade, rastreabilidade e revisão humana}
\label{rt-2-5-3}

% V5: a introdução final anuncia explicabilidade, rastreabilidade e revisão humana no Cap. 2. Phillips et al. foi conferido (NISTIR 8312, seção 2, p. 2-5).
Um parecer só pode ser revisado se for possível entender como ele foi produzido. \citeonline{phillips2021} propõem quatro princípios para sistemas de inteligência artificial explicáveis. O sistema deve fornecer evidências ou razões para suas saídas, e essas explicações devem ser compreensíveis para quem as recebe e refletir corretamente o processo que gerou a saída. Além disso, o sistema deve operar apenas nas condições para as quais foi projetado. Os autores distinguem a exatidão da explicação da exatidão da decisão, de modo que uma explicação fiel não torna a decisão correta.

A rastreabilidade sustenta essa explicação. A proveniência é a informação sobre as entidades, as atividades e as pessoas envolvidas na produção de um dado, e pode ser usada para avaliar sua qualidade, confiabilidade ou credibilidade \cite{groth2013}. A organização de tarefas, a vinculação às fontes e o registro da execução também são discutidos em propostas de fluxos governáveis e auditáveis \cite{rousek2026}.

Os registros permitem inspecionar o procedimento, mas não asseguram a correção de cada interpretação. Fluxos apoiados por modelos de linguagem podem produzir afirmações sem suporte, e por isso a ligação entre evidência, processamento e resultado precisa ser preservada \cite{rousek2026}. \citeonline{sowan2026} reconhecem limite semelhante ao afirmar que a utilidade de suas trilhas de decisão para conformidade e responsabilização ainda depende de validação. Daí a importância da revisão humana. \citeonline{carvalho2026} defendem que a aceitação de riscos permaneça com uma pessoa ou instância identificável.

Neste trabalho, a explicabilidade é tratada como rastreabilidade estrutural. O parecer aponta as evidências, a origem de cada uma e a regra aplicada, preservando os registros que sustentam essa ligação. A utilidade dessas explicações para as pessoas não decorre dos registros. Ela é examinada pela apreciação dos especialistas, definida no Capítulo~\ref{cap-metodologia}.

\section{Trabalhos Relacionados}
\label{rt-2-6}

% V6: seção reescrita. Cada trabalho é descrito pelo problema, pelo método, pelos agentes e sua comunicação, pela consolidação, pela rastreabilidade e pela avaliação relatada.
% As descrições foram conferidas nos PDFs locais. Só se afirma a ausência de um elemento quando o próprio texto do trabalho a sustenta.
As seções anteriores apresentaram os conceitos que sustentam a proposta. Esta seção examina nove trabalhos selecionados pela pesquisa bibliográfica descrita na Seção ??. Cada um é descrito pelo problema que trata, pela forma como organiza agentes e comunicação, pela maneira de consolidar resultados, pelos registros que conserva e pela avaliação que relata. A proximidade temática não significa equivalência de domínio, de método ou de resultado.

\subsection{Apoio à decisão sobre evidências de software}
\label{rt-2-6-1}

\citeonline{zuev2025} tratam da decisão de incorporar código ao ramo principal de um projeto. A abordagem reúne os relatórios de análises estática e dinâmica em um módulo de apoio à decisão, que calcula indicadores e aplica regras para apresentar uma recomendação ao usuário. O experimento relatado usa uma amostra de cerca de 100 varreduras e ilustra o funcionamento do módulo. O artigo não descreve uma arquitetura de agentes, e a seção experimental não relata o julgamento de especialistas sobre as recomendações. Para esta pesquisa, interessa a passagem dos resultados das ferramentas para uma recomendação apoiada em regras.

\citeonline{khoee2026} apresentam o GateLens, voltado à análise de dados de \textit{releases} de software automotivo. A solução usa álgebra relacional como representação intermediária entre a consulta em linguagem natural e o código executável, o que permite inspecionar cada transformação. Os autores contrastam essa escolha com sistemas multiagentes e de planejamento, e a solução não depende de orquestração entre agentes. A avaliação compara as respostas a consultas com um gabarito preparado por especialistas. Para esta pesquisa, interessa a possibilidade de verificar o caminho entre a entrada e o resultado.

\subsection{Comunicação e consolidação por agentes especializados}
\label{rt-2-6-2}

\citeonline{lazarin2025} tratam da interação entre pessoas e agentes embarcados. Os autores desenvolvem uma arquitetura de agente baseada em Jason \cite{bordini2007}, na qual mensagens KQML trafegam por protocolos de correio eletrônico. O trabalho demonstra a arquitetura em provas de conceito. Interessa à pesquisa a separação entre a intenção comunicativa da mensagem e o meio de transporte. O domínio é distinto da auditoria de repositórios.

\citeonline{yesil2026} tratam do controle de acesso em recuperação aumentada por geração multimodal, técnica em que o modelo responde com apoio em documentos recuperados. Um agente orquestrador coordena o fluxo e não trata da autorização, que fica com agentes especializados em verificações textuais e visuais. Os agentes se comunicam por atos de fala, e um agente de decisão combina os resultados por disjunção lógica, isto é, o acesso é autorizado se uma das verificações o permitir. A lógica é expressa em planos que podem ser inspecionados. A avaliação usa 167 pôsteres de seis instituições. Interessa à pesquisa a separação entre coordenação, verificação especializada e consolidação. A regra de combinação, porém, responde ao problema de autorização, e o fluxo descrito não inclui a revisão das conclusões dos especialistas a partir de consultas entre eles.

\subsection{Políticas e fluxos governáveis}
\label{rt-2-6-3}

\citeonline{ait2025} tratam da governança de projetos de software em que pessoas e agentes colaboram. Os autores propõem uma linguagem específica de domínio para representar participantes, papéis, escopos e políticas, com um mecanismo de decisão que aplica essas políticas, e apresentam uma prova de conceito. A avaliação da proposta é indicada como trabalho futuro. Interessa à pesquisa a forma como uma regra é ligada aos participantes e ao escopo em que vale.

\citeonline{yan2025} tratam da aplicação de normas e de sanções do ponto de vista do agente. Os autores estendem uma linguagem de programação normativa, na qual normas e sanções são abstrações de primeira classe, e propõem uma arquitetura de agente BDI normativo. A proposta é aplicada a um estudo de caso, e a experimentação é deixada para trabalhos futuros. A distinção entre processar uma norma e aplicar uma sanção ajuda a delimitar um agente que emite um parecer e não executa consequências.

\citeonline{rousek2026} apresenta um \textit{framework} conceitual para fluxos governáveis de apoio à decisão em defesa. O \textit{framework} tem três camadas, que são a decomposição explícita de tarefas, a vinculação das afirmações às evidências, com proveniência, e o controle delimitado, com registro e pontos de autorização. O autor declara que o estudo não valida o \textit{framework} empiricamente. A obra é usada aqui como referência conceitual para os registros e para a delimitação da autonomia.

\subsection{Separação de funções e registros de decisão}
\label{rt-2-6-4}

\citeonline{mosciatti2026} tratam de operações de segurança em dispositivos e aplicações com interfaces web. A arquitetura divide a tarefa em fases, e cada fase fica com um agente que tem memória, permissões e ferramentas isoladas. Os autores avaliam um protótipo de rotação de senhas em 12 aplicações e relatam sucesso em todas as execuções do experimento. Interessa à pesquisa a separação de deveres entre agentes e a existência de uma fase própria de validação. O trabalho executa operações de segurança, e não emite pareceres sobre repositórios.

\citeonline{sowan2026} apresentam o AEGIS-Chain, uma arquitetura de apoio à decisão em cadeias de suprimentos. A arquitetura combina limites operacionais, níveis de escalonamento para pessoas e uma trilha de auditoria das decisões. A avaliação usa mil episódios de simulação. Os autores reconhecem que a utilidade dos registros para conformidade e responsabilização não foi validada empiricamente. Interessa à pesquisa a identificação das situações que pedem encaminhamento a uma pessoa.

\subsection{Comparação e posicionamento da proposta}
\label{rt-2-6-5}

A Tabela~\ref{quadro-trabalhos-relacionados} compara os trabalhos em quatro aspectos. As células descrevem o que cada texto relata.

\begin{table}[htbp]
\centering
\caption{Comparação dos trabalhos relacionados}
\label{quadro-trabalhos-relacionados}
\begingroup
\scriptsize
\SingleSpacing
\setlength{\tabcolsep}{3pt}
\renewcommand{\arraystretch}{1.15}
\begin{tabularx}{\textwidth}{>{\raggedright\arraybackslash}p{2.1cm}>{\raggedright\arraybackslash}X>{\raggedright\arraybackslash}X>{\raggedright\arraybackslash}X>{\raggedright\arraybackslash}X}
\hline
\textbf{Trabalho} & \textbf{Agentes e comunicação} & \textbf{Consolidação} & \textbf{Rastreabilidade} & \textbf{Avaliação relatada} \\ \hline
\citeonline{zuev2025} & Módulo único, sem arquitetura de agentes & Indicadores e regras sobre os relatórios das análises & Relatório combinado das análises & Amostra de cerca de 100 varreduras \\ \hline
\citeonline{khoee2026} & Agente único, sem orquestração entre agentes & Não se aplica & Representação intermediária em álgebra relacional & Consultas comparadas com gabarito de especialistas \\ \hline
\citeonline{lazarin2025} & Agentes BDI e pessoas, com mensagens KQML por correio eletrônico & Não é o foco do trabalho & Não é o foco do trabalho & Provas de conceito \\ \hline
\citeonline{yesil2026} & Orquestrador e agentes especializados, com atos de fala & Agente de decisão, por disjunção lógica & Planos inspecionáveis & 167 pôsteres de seis instituições \\ \hline
\citeonline{ait2025} & Pessoas e agentes como participantes de políticas & Mecanismo de decisão que aplica as políticas & Políticas explícitas, com papéis e escopos & Prova de conceito, com avaliação futura \\ \hline
\citeonline{yan2025} & Agentes BDI normativos & Processamento de normas e de sanções pelo agente & Normas e sanções explícitas & Estudo de caso, com experimentação futura \\ \hline
\citeonline{rousek2026} & Cinco papéis inspecionáveis & Controle delimitado, com pontos de autorização & Vinculação às evidências, com proveniência & Sem validação empírica, declarada pelo autor \\ \hline
\citeonline{mosciatti2026} & Um agente por fase, com memória e permissões isoladas & Fase própria de validação & Não é o foco do trabalho & Rotação de senhas em 12 aplicações \\ \hline
\citeonline{sowan2026} & Arquitetura de uma única organização & Limites operacionais e escalonamento a pessoas & Trilha de auditoria das decisões & Mil episódios de simulação \\ \hline
\end{tabularx}
\endgroup
\legend{Fonte: elaboração do autor, a partir da leitura dos trabalhos.}
\end{table}

A comparação mostra três pontos. O primeiro é que a passagem dos resultados de ferramentas para uma recomendação por regras já aparece em \citeonline{zuev2025}, sem uma arquitetura de agentes. O segundo é que a separação entre coordenação, verificação especializada e consolidação aparece em \citeonline{yesil2026} e em \citeonline{mosciatti2026}, em domínios diferentes da auditoria de repositórios. O terceiro é que, nos trabalhos lidos que têm uma etapa de consolidação, os textos descrevem a entrega dos resultados a quem consolida. Eles não descrevem a consulta direta entre os especialistas para rever uma conclusão antes disso.

Esse terceiro ponto delimita o que esta proposta examina. Ela reúne, na auditoria de repositórios de automações, a consulta direta entre agentes especialistas, o contexto documental dos componentes e a consolidação por regras em um parecer destinado a um responsável humano. A afirmação se restringe aos nove trabalhos lidos. Ela não significa que a literatura não contenha soluções equivalentes, e a proposta não pressupõe eficácia demonstrada.

\section{Resumo do Capítulo}
\label{rt-resumo}

A auditoria de repositório avalia uma versão identificada do código e de seus artefatos. Suas atividades produzem achados de qualidade, de segurança e de dependências, e as evidências permitem conferi-los. O contexto documental indica o que cada achado significa para a organização. A conclusão do especialista registra a interpretação de cada área, e o parecer reúne essas conclusões para apoiar a decisão do responsável humano. Esse apoio não constitui uma auditoria independente no sentido das normas.

Os sistemas multiagentes distribuem essas atividades entre especialistas, que trocam informações por comunicação P2P. Um modelo de linguagem apoia apenas a obtenção do contexto documental. O Agente Auditor, que exerce o controle semi-centralizado, reúne as conclusões dos especialistas e as consolida por regras derivadas de insumos de governança. Os registros de proveniência permitem explicar e revisar como cada parecer foi produzido. As tecnologias que materializam esses elementos são descritas no Capítulo~\ref{cap-suporte-tecnologico}, os procedimentos de coleta e análise, no Capítulo~\ref{cap-metodologia}, e os papéis, os fluxos e as regras da solução, no Capítulo~\ref{cap-proposta}.